\chapter{METHODOLOGY}
\section{Study Design}
The assessment follows the sequence of model definition, support-spring representation, modal comparison, response-spectrum analysis, gravity-load response examination and fragility interpretation. This ordering relates each research question to its corresponding engineering demand parameter and preserves a consistent baseline for comparing the four support cases.

The six structural blocks of the Infectious and Communicable Disease Hospital were analysed for fixed-base, hard-soil, medium-soil and soft-soil support conditions. Modal, response-spectrum and gravity-load results were compared to evaluate the influence of foundation flexibility on natural period, displacement, inter-storey drift, shear, settlement and member forces. Fragility relationships were considered as a complementary measure of seismic performance.

\begin{figure}[htbp]
\centering
\includegraphics[width=0.92\textwidth,height=0.60\textheight,keepaspectratio]{figures/ch3/analysis_flowchart_seq057.pdf}
\caption{Research methodology for evaluating soil--structure interaction and seismic fragility across the six hospital blocks}
\label{fig:research_method_flowchart}
\end{figure}
\FloatBarrier

\section{Hospital and Structural Configuration}
The case study is the existing \textbf{Infectious and Communicable Disease Hospital} in Mohoriya, Lekhnath, Pokhara Metropolitan City--30, Kaski, Gandaki Province. The hospital is also known under its later institutional name, Provincial General and Communicable Disease Hospital. The submitted architectural plan identifies six structural blocks modelled as G+4 (ground plus four upper floors).\footnote{Provincial hospital institutional website and contact details: \url{https://icdh.gandaki.gov.np/} and \url{https://icdh.gandaki.gov.np/contact}.}

\begin{figure}[htbp]\centering
\includegraphics[width=0.90\textwidth,page=1]{seq050_assets/submitted_plan.pdf}
\caption{Architectural plan of the G+4 hospital, showing the six submitted blocks}
\label{fig:hospital_plan_methodology}
\end{figure}

\begin{figure}[htbp]\centering
\includegraphics[width=0.74\textwidth]{seq050_assets/b1_new.png}
\caption{Three-dimensional structural representation of Block 1; views for all six blocks appear in the appendix}
\label{fig:hospital_3d_methodology}
\end{figure}

Each block was evaluated as a separate structural model, consistent with the building plan. The member dimensions entered for each block are summarized in Table~\ref{tab:member_dimensions_methodology}. The blocks were compared individually to account for differences in their structural configurations.

\begingroup
\normalsize
\renewcommand{\arraystretch}{1.15}
\begin{longtable}{@{}p{1.15cm}p{3.15cm}>{\raggedright\arraybackslash}p{9.0cm}@{}}
\caption{Principal structural member dimensions of the six hospital blocks}\label{tab:member_dimensions_methodology}\\
\toprule
Block & Member type & Dimensions (mm) \\
\midrule
\endfirsthead
\multicolumn{3}{l}{\normalsize Table 3.1 continued}\\
\toprule
Block & Member type & Dimensions (mm) \\
\midrule
\endhead
\bottomrule
\endlastfoot
1 & Columns & 700$\times$700, 750$\times$750 \\
 & Beams & 350$\times$550; 400$\times$800; 350$\times$700; 300$\times$500; 350$\times$650; 250$\times$550; 250$\times$350 \\
 & Wall thickness & 200 \\
 & Slab thickness & 150 \\
\addlinespace[0.35em]
2 & Columns & 700$\times$700, 550$\times$550 \\
 & Beams & 400$\times$700; 350$\times$650; 300$\times$500; 500$\times$800; 300$\times$450; 350$\times$500; 250$\times$350 \\
 & Wall thickness & -- \\
 & Slab thickness & 150 \\
\addlinespace[0.35em]
3 & Columns & 700$\times$700 \\
 & Beams & 350$\times$600; 400$\times$800; 400$\times$750; 350$\times$550; 300$\times$500; 250$\times$500; 250$\times$350 \\
 & Wall thickness & -- \\
 & Slab thickness & 150 \\
\addlinespace[0.35em]
4 & Columns & 600$\times$600, 650$\times$650 \\
 & Beams & 350$\times$600; 350$\times$650; 400$\times$750; 250$\times$450; 250$\times$350 \\
 & Wall thickness & -- \\
 & Slab thickness & 150 \\
\addlinespace[0.35em]
5 & Columns & 700$\times$700 \\
 & Beams & 350$\times$550; 400$\times$750; 300$\times$500; 350$\times$650; 350$\times$500; 300$\times$450; 250$\times$350 \\
 & Wall thickness & 200 \\
 & Slab thickness & 150 \\
\addlinespace[0.35em]
6 & Columns & 700$\times$700, 550$\times$550 \\
 & Beams & 350$\times$550; 400$\times$750; 300$\times$450; 250$\times$350; 350$\times$450 \\
 & Wall thickness & 200 \\
 & Slab thickness & 150 \\
\addlinespace[0.35em]
\end{longtable}
\noindent\textit{Note:} ``--'' = wall not specified. The staircase slabs in Blocks 1 and 6 and the ramp slab in Block 2 are 175~mm thick.
\endgroup



\section{Numerical Model}
The three-dimensional reinforced-concrete numerical models of the hospital blocks were developed using ETABS v22, a finite-element-based structural analysis program. The six blocks shown in the architectural plan were represented separately. Structural members include beams, columns and slabs, with shear-wall elements where indicated in the submitted structural dimensions. Member sections were assigned according to Table~\ref{tab:member_dimensions_methodology}; the foundation-support properties follow Table~\ref{tab:spring_stiffness_methodology}. Building geometry, imposed gravity loads and seismic loading definitions were taken from the corresponding researcher-supplied model schedules.

For each block, the fixed-base idealization restrains foundation movement and provides the reference response. Flexible-base cases represent hard, medium and soft soil using directional translational and rotational foundation springs. In generalized coordinates, the foundation stiffness may be expressed as
\begin{equation}
\mathbf{K}_f=\operatorname{diag}(K_x,K_y,K_z,K_{\theta x},K_{\theta y},K_t),
\label{eq:foundation_six_dof}
\end{equation}
\RegisterThesisEquation{Six-degree-of-freedom foundation spring matrix}
where the first three terms represent translational restraint and the final three represent rocking and torsion about the relevant local axes. This diagonal description is the idealized uncoupled six-degree-of-freedom spring model; it should not be interpreted as a measured full impedance matrix for the as-built foundation.

The Winkler idealization relates soil reaction locally to displacement, whereas the equivalent directional foundation impedances are informed by the Gazetas (1991) elastic-foundation framework \cite{wolf1985,gazetas1991}. The reported input schedules provide the actual six spring-stiffness components for each block and soil condition. The specified stiffness coefficients define the hard-, medium- and soft-soil analytical cases; these models idealize the supporting soil through linear elastic springs. Figure~\ref{fig:winkler_gazetas_schematic} illustrates the mechanical distinction.

\begin{figure}[htbp]
\centering
\includegraphics[width=0.90\textwidth]{figures/ch2/winkler_gazetas_schematic.pdf}
\caption{Conceptual Winkler springs and six-component equivalent foundation restraint for flexible-base structural modelling (schematic; not to scale).}
\label{fig:winkler_gazetas_schematic}
\end{figure}

\subsection{Analysis cases and response quantities}
Modal analysis provides natural periods and cumulative participating mass ratios in the two horizontal directions. Response-spectrum analysis gives the horizontal displacement, inter-storey drift, base reaction and storey-shear response for the specified X and Y cases. The independent $1.2DL+1.5LL$ gravity combination provides base-joint settlement, vertical reaction and beam/column member-force components. Each set of results is tabulated against its own load-case definition to prevent the erroneous use of gravity results as earthquake results. The X-labelled fragility workbooks provide source curves for comparative interpretation; new probabilistic fits are not inferred from the response-spectrum maxima alone.

\subsection{Comparative assessment procedure}
The response quantities are evaluated separately for each structural block and its prescribed foundation-support condition. Trends are examined using the corresponding load cases, units, component directions and modal or seismic response definitions. Comparing the flexible-support cases with the fixed-base case identifies the variation associated with the adopted foundation stiffness within each block.


\section{Seismic Parameters and Gravity Loading}
The structural models adopt an importance factor of 1.5, a damping ratio of 5 percent, and a Special Moment Resisting Frame (SMRF) structural-system designation. Seismic loading and site classification follow NBC 105:2025 \cite{nbc1052025}. The gravity-load assignments for the model are summarized in Table~\ref{tab:gravity_load_inputs}. The imposed floor loading differs according to the functional use of the spaces and the relevant structural element.

\begin{table}[htbp]
\centering
\caption{Gravity load inputs assigned in the structural analysis}
\label{tab:gravity_load_inputs}
\renewcommand{\arraystretch}{1.12}
\begin{tabular}{@{}p{5.5cm}p{3cm}p{3.3cm}@{}}
\toprule
Load item & Assigned load & Application \\
\midrule
Floor finish & $2.5\ \mathrm{kN/m^{2}}$ & Slab \\
Wall load & $16.5\ \mathrm{kN/m}$ & Beam \\
Ramp parapet wall & $5\ \mathrm{kN/m}$ & Ramp beam \\
Water tank load & $3.6\ \mathrm{kN/m^{2}}$ & Fifth-floor slab \\
Ward floor live load & $4\ \mathrm{kN/m^{2}}$ & Slab \\
Corridor live load & $4\ \mathrm{kN/m^{2}}$ & Slab \\
Operation theatre and ICU live load & $4\ \mathrm{kN/m^{2}}$ & Slab \\
Staircase and ramp floor live load & $5\ \mathrm{kN/m^{2}}$ & Stair/ramp slab \\
Roof live load & $0.75\ \mathrm{kN/m^{2}}$ & Roof slab \\
\bottomrule
\end{tabular}
\end{table}
\FloatBarrier

Response-spectrum analysis provides the linear seismic demand comparison for the fixed- and flexible-support models. Modal participating mass is examined in the two principal horizontal directions, with the NBC 105:2025 cumulative participating-mass criterion of at least 90 percent used for assessment \cite{nbc1052025}. Gravity-load combinations and spectral-response cases are interpreted independently when comparing foundation and superstructure demand.

\section{Soil Parameters and SSI Idealization}
The study considers three flexible-base soil conditions, designated hard, medium, and soft. Their site-class categories and material parameters used for the analytical comparison are summarized in Table~\ref{tab:soil_property_inputs}. The hard-soil support corresponds to the greatest subgrade stiffness and the soft-soil case to the smallest; these properties establish the comparative foundation spring idealizations.

\begin{table}[htbp]
\centering
\caption{Soil parameters and site classifications used in SSI modelling}
\label{tab:soil_property_inputs}
\fontsize{9}{11}\selectfont
\setlength{\tabcolsep}{3.5pt}
\begin{tabular}{@{}lccrrrrr@{}}
\toprule
Soil case & NBC class & NBC $V_{s,30}$ (m/s) & Project $V_{s,30}$ (m/s) & $E$ (MPa) & $\nu$ & $G$ (MPa) & SBC (kPa) \\
\midrule
Hard & B & $350<V_{s,30}\leq800$ & 575 & 1550 & 0.30 & 595 & 400 \\
Medium & C & $150<V_{s,30}\leq350$ & 250 & 300 & 0.35 & 113 & 200 \\
Soft & D & $V_{s,30}\leq150$ & 125 & 80 & 0.40 & 28 & 100 \\
\bottomrule
\end{tabular}
\end{table}
\FloatBarrier

Within the Winkler idealization, the contact pressure at a foundation point is related to local vertical soil deformation through a modulus of subgrade reaction. Its one-parameter vertical form is
\begin{equation}
p(x,y)=k_s\,w(x,y),
\label{eq:winkler_soil_reaction}
\end{equation}
\RegisterThesisEquation{Local vertical soil reaction in the Winkler idealization}
where $p(x,y)$ is the foundation reaction per unit area, $w(x,y)$ denotes the local foundation displacement, and $k_s$ is the modulus of subgrade reaction. The six-component foundation-spring representation extends beyond this scalar vertical relationship to independent horizontal, vertical, rocking, and torsional restraints. The equivalent static stiffness components are formulated using elastic half-space foundation relations as discussed by Gazetas \cite{gazetas1991}. For the square isolated footing adopted in the spring derivation, the plan dimensions are $3.0\,\mathrm{m}\times3.0\,\mathrm{m}$, with $L=B=1.5\,\mathrm{m}$ in the $2L\times2B$ notation. The resulting foundation area is $9\,\mathrm{m^2}$ and the geometric properties are $I_{bx}=I_{by}=6.75\,\mathrm{m^4}$ and $J_b=13.5\,\mathrm{m^4}$.

\section{Foundation-Support Representation}
Soil--structure interaction is represented by springs describing translational, rocking and torsional foundation restraint. The fixed-base case provides the comparison reference; the flexible-base cases are categorized as hard, medium and soft soil. The discrete-spring concept is compatible with substructure SSI modelling \cite{wolf1985,gazetas1991}. The model incorporates the translational and rotational spring stiffnesses assigned to the foundation for each soil condition.

The equivalent static foundation-stiffness expressions used for comparison follow the elastic-half-space formulation of Gazetas (1991) \cite{gazetas1991}. For the present square footing with plan dimensions $2L\times2B$ and $L=B$, define $\chi=A_b/(4L^2)=B/L=1$. The vertical, horizontal, rocking and torsional expressions are
\begin{equation}
K_z = \frac{2GL}{1-\nu}\left(0.73+1.54\chi^{0.75}\right)\label{eq:gazetas_kz}
\end{equation}
\RegisterThesisEquation{Gazetas vertical foundation stiffness}
\begin{equation}
K_x=K_y = \frac{2GL}{2-\nu}\left(2+2.50\chi^{0.85}\right)\label{eq:gazetas_kxy}
\end{equation}
\RegisterThesisEquation{Gazetas horizontal foundation stiffness}
\begin{equation}
K_{\theta x} = \frac{G}{1-\nu}I_{bx}^{0.75}\left(\frac{L}{B}\right)^{0.25}\left(2.4+0.5\frac{B}{L}\right)\label{eq:gazetas_krx}
\end{equation}
\RegisterThesisEquation{Gazetas rocking stiffness about the x-axis}
\begin{equation}
K_{\theta y} = \frac{G}{1-\nu}I_{by}^{0.75}\left[3\left(\frac{L}{B}\right)^{0.15}\right]\label{eq:gazetas_kry}
\end{equation}
\RegisterThesisEquation{Gazetas rocking stiffness about the y-axis}
\begin{equation}
K_t = GJ_b^{0.75}\left[4+11\left(1-\frac{B}{L}\right)^{0.75}\right]\label{eq:gazetas_kt}
\end{equation}
\RegisterThesisEquation{Gazetas torsional foundation stiffness}
Here, $G$ denotes the soil shear modulus, $\nu$ is Poisson's ratio, $I_{bx}$ and $I_{by}$ are the second moments of the foundation contact area, and $J_b=I_{bx}+I_{by}$ is its polar area moment. For the square $3.0\,\mathrm{m}\times3.0\,\mathrm{m}$ footing, $L=B=1.5\,\mathrm{m}$, $A_b=9\,\mathrm{m^2}$, $I_{bx}=I_{by}=6.75\,\mathrm{m^4}$ and $J_b=13.5\,\mathrm{m^4}$. The tabulated values remain the original foundation-spring inputs assigned in ETABS; the equations describe the analytical stiffness relationships without replacing the recorded model inputs.

The spring-stiffness matrix includes six components ($K_x$, $K_y$, $K_z$, $K_{\theta x}$, $K_{\theta y}$ and $K_t$) for each of three flexible-support conditions in every block. Table~\ref{tab:spring_stiffness_methodology} presents all 18 block/soil rows, for the six structural blocks.

\begin{table}[!htbp]
\caption{Foundation spring input values by block and soil condition (Table 3.4)}
\label{tab:spring_stiffness_methodology}\centering\tiny
\resizebox{\textwidth}{!}{%
\begin{tabular}{clrrrrrr}
\hline\hline
Block & Soil & $K_x$ & $K_y$ & $K_z$ & $K_{\theta x}$ & $K_{\theta y}$ & $K_t$ \\ \hline
1 & Hard & 3392594.941 & 3392594.941 & 6753250.000 & 16392119.903 & 16957365.417 & 38025034.110 \\
1 & Medium & 663832.430 & 663832.430 & 1381207.692 & 3352596.469 & 3468203.244 & 7777065.800 \\
1 & Soft & 169629.747 & 169629.747 & 370766.667 & 899959.524 & 930992.611 & 2087648.932 \\
2 & Hard & 3392594.941 & 3392594.941 & 6753250.000 & 16392119.903 & 16957365.417 & 38025034.110 \\
2 & Medium & 663832.430 & 663832.430 & 1381207.692 & 3352596.469 & 3468203.244 & 7777065.800 \\
2 & Soft & 169629.747 & 169629.747 & 370766.667 & 899959.524 & 930992.611 & 2087648.932 \\
3 & Hard & 4119579.600 & 4119579.600 & 9004333.333 & 28546821.100 & 29531194.200 & 66220467.570 \\
3 & Medium & 758665.634 & 758665.634 & 1578523.077 & 5004458.869 & 5177026.416 & 11608914.547 \\
3 & Soft & 193862.568 & 193862.568 & 423733.333 & 1343379.814 & 1389703.256 & 3116257.297 \\
4 & Hard & 4119579.600 & 4119579.600 & 9004333.333 & 28546821.100 & 29531194.200 & 66220467.570 \\
4 & Medium & 758665.634 & 758665.634 & 1578523.077 & 5004458.869 & 5177026.416 & 11608914.547 \\
4 & Soft & 193862.568 & 193862.568 & 423733.333 & 1343379.814 & 1389703.256 & 3116257.297 \\
5 & Hard & 2907938.521 & 2907938.521 & 5788500.000 & 10322734.400 & 10678690.758 & 23945794.075 \\
5 & Medium & 568999.226 & 568999.226 & 1183892.308 & 2111256.085 & 2184058.019 & 4897510.824 \\
5 & Soft & 145396.926 & 145396.926 & 317800.000 & 566738.359 & 586281.061 & 1314671.047 \\
6 & Hard & 2907938.521 & 2907938.521 & 5788500.000 & 10322734.400 & 10678690.758 & 23945794.075 \\
6 & Medium & 568999.226 & 568999.226 & 1183892.308 & 2111256.085 & 2184058.019 & 4897510.824 \\
6 & Soft & 145396.926 & 145396.926 & 317800.000 & 566738.359 & 586281.061 & 1314671.047 \\
\hline\hline
\end{tabular}}
\par\smallskip\footnotesize The numerical values represent the translational and rotational spring stiffnesses used in the foundation models.
\end{table}
\FloatBarrier

\section{Modal Response Extraction and Comparison}
Modal analysis results were evaluated separately for each block and foundation-support condition. The first-mode period $T_{1,b,s}$ is extracted using the row and column headings of the actual worksheet. For every case, the smallest modal number with both horizontal cumulative participating-mass ratios at or above 0.90 is recorded. 

The percentage change in period from fixed to a flexible support case is computed within the same block as
\begin{equation}
\Delta T_{b,s}=100\frac{T_{1,b,s}-T_{1,b,\mathrm{fixed}}}{T_{1,b,\mathrm{fixed}}}\,\%,
\label{eq:ssi_sensitivity}
\end{equation}
\RegisterThesisEquation{Percentage change in fundamental period}
where $b$ identifies the block and $s$ denotes hard, medium or soft support. No percentage comparison is calculated between unrelated blocks or between physically different output quantities.

\section{Response Charts and Drift Screening}
Displacement, drift and fragility relationships were examined for the different foundation-support conditions. The displacement and drift responses were obtained from the corresponding directional analyses, while fragility relationships were examined separately for each structural block. The relevant results are presented graphically in Chapter~4 and the appendices.

The adopted inter-storey drift ratio (ISDR) performance classification is shown in Table~\ref{tab:legacy_isdr_performance_bands}. Immediate Occupancy (IO), Life Safety (LS) and Collapse Prevention (CP) are standard performance concepts \cite{fema356}, whereas the numerical percentages below are the study's adopted ISDR screening intervals.

\begin{table}[htbp]
\caption{Study-adopted ISDR performance bands for fragility screening (Table 3.5)}
\label{tab:legacy_isdr_performance_bands}\centering
\begin{tabular}{lc}
\hline\hline
Performance level & Maximum ISDR \\ \hline
Immediate Occupancy (IO) & $\leq 1.0\%$ \\
Life Safety (LS) & $>1.0\%$ to $\leq 2.0\%$ \\
Collapse Prevention (CP) & $>2.0\%$ to $\leq 4.0\%$ \\
Collapse / Failure (beyond CP) & $>4.0\%$ \\ \hline\hline
\end{tabular}
\end{table}

The ISDR bands are descriptive demand categories. The adopted screening limits are used to interpret drift demand; fragility probabilities are interpreted according to the damage-state definitions associated with their respective curves.

\section{Interpretation Procedure}
Results are organized by block number, support case, output quantity, and original load case. For the response-spectrum displacement results, the reported response at the greatest listed storey elevation is extracted separately from the RS-X X-component and RS-Y Y-component tables. For response-spectrum drift, the maximum absolute reported storey drift ratio is identified separately for the same response cases and converted to percent. Base shear is read from the response-spectrum maximum FX (RS-X) and FY (RS-Y) base-reaction rows; the largest absolute response-spectrum storey-shear value is selected from the corresponding storey rows without summation across floors.

Foundation settlements are the largest absolute base-joint Uz displacement, vertical reactions are the FZ component, and beam and column force envelopes are maximum absolute member-force components at the source-recorded element stations, all under the independent \texttt{1.2DL + 1.5LL} combination. These gravity-combination outputs are not relabelled as response-spectrum results. The controlling member forces were evaluated for each block and support condition. Results from inconsistent model assignments were excluded from direct quantitative comparison.
